\begin{musica}{Ondas do Douro}{Tuna Universitária do Porto}{}
\nota{Introdução: D A D D7 G Gm D A D}

\nota{Refrão:}
Linda don\ac{D}{ze}la vem à janela que a Tuna \ac{A}{pa}ssa,
Ouve este canto que o teu encanto enche de \ac{D}{gra}ça.
Olha p’rá Lua que a noite é \ac{D7}{tua} e o tro\ac{G}{va}d\ac{Gm}{or}
Enamo\ac{D}{ra}do, canta ele\ac{A}{va}do trovas da a\ac{D}{mor}. \ac{F#}{}__

São teus ca\ac{Bm}{be}los ondas que o Douro leva p’ró \ac{F#}{mar},
Lento embalo de melodia que faz so\ac{Bm}{nhar},
Barcos ra\ac{B7}{be}los feitos de esperança de um teu o\ac{Em}{lhar}
E a Tuna ron\ac{Bm}{da} junto à Ri\ac{F#}{bei}ra p’ra te can\ac{B7}{tar}. \ac{A}{}_

\nota{Refrão}

Levo nos olhos a tua imagem, brando fulgor,
Levo a Saudade, deixo esta trova ao teu amor;
Põe um sorriso, não te entristeças se a Tuna parte.
Que o estudante é terno amante, virá cantar-te.

\posicoes{D: xx0232, D7: xx0212, Em: 022000, F#: 244322, G: 320033, Gm: 355333,
  A: x02220, B7: 021202, Bm: x24432}
\end{musica}
